% This file contains the abstract of the thesis

Proiectul VisioScience3D vine ca un răspuns la nevoia de a crea un mediu de învățare interactiv
și captivant pentru elevii din învățământul preuniversitar ce lipsește de pe piața curentă.
Acesta îmbină tehnologii moderne pentru a crea un sistem de învățare bazat pe vizualizări 3D
și simulări interactive, care să faciliteze înțelegerea conceptelor complexe din domeniul științelor exacte.

În momentul curent învățarea geometriei, fizicii, chimiei și altor discipline științifice
se face prin metode tradiționale, care nu reușesc să capteze atenția elevilor.
Prin acest proiect se dorește venirea în întâmpinarea acestei nevoi și oferirea
unui mediu de învățare interactiv și captivant care să faciliteze
activitatea didactică și să îmbunătățească rezultatele elevilor.

Proiectul VisioScience3D este o aplicație web care permite utilizatorilor să exploreze
domeniul științelor exacte prin intermediul simulărilor interactive și vizualizărilor 3D.
Acesta oferă o soluție inovatoare și complexă pentru elevi cât în egală masură și pentru
profesori, având posibilitatea să predea și să evalueze elevii prin intermediul acestuia rapid și eficient.